% Syntra LaTeX Template
% For Arabic support, use XeLaTeX with:
%   \usepackage{fontspec}
%   \usepackage{polyglossia}
%   \setdefaultlanguage{english}
%   \setotherlanguage{arabic}
%   \newfontfamily\arabicfont[Script=Arabic]{Amiri}

\documentclass[12pt,a4paper]{article}

% Font and encoding
\usepackage[utf8]{inputenc}
\usepackage[T1]{fontenc}

% Math packages
\usepackage{amsmath,amssymb}

% Table support
\usepackage{longtable,booktabs}

% Graphics and links
\usepackage{graphicx}
\usepackage{hyperref}
\usepackage{geometry}

% Page geometry
\geometry{
    left=2.5cm,
    right=2.5cm,
    top=3cm,
    bottom=3cm
}

% Fix pandoc tightlist issue
\providecommand{\tightlist}{\setlength{\itemsep}{0pt}\setlength{\parskip}{0pt}}

% Hyperref settings
\hypersetup{
    colorlinks=true,
    linkcolor=blue,
    filecolor=magenta,
    urlcolor=cyan,
    pdftitle={Syntra Research Documentation},
    pdfauthor={Syntra-Env},
    pdfpagemode=FullScreen
}

\begin{document}

\section{Text Substitution as a Design Pattern for Scope Control and
Cognitive Load
Optimization}\label{text-substitution-as-a-design-pattern-for-scope-control-and-cognitive-load-optimization}

\includegraphics[width=7.29167in,height=5.20833in,alt={Text Substitution Interface}]{https://github.com/user-attachments/assets/7d23a9f0-4e79-448e-81fa-e9779fc6543f}

In line with the cognitive load section in page about write-read
juxtaposition:

There\textquotesingle s this bar at the bottom with different buttons
that represent different types of substitutions you can perform, maybe
you want to see the details of a section of the text that are relevant
to syntax, or morphology, etc.

So, what if you didn\textquotesingle t have to open any other panel or
page at all. What if all we did is add and subtract context, and
provided tools for doing this (the buttons means for different kinds of
substitution.)

The set of actions would be:

\begin{itemize}
\tightlist
\item
  Click on a button to select the type of substitution
\item
  Left click on a portion of the text to use substitutive expansion and
  right click for contraction
\item
  Maybe add another dimension of navigability by using the mouse scroll.
  Not in a way that would actually scroll up and down but maybe some
  other action that has to do with going up or down the layers of
  substitution.
\end{itemize}

Part of the vision was the mouse icon itself changing color or style or
whatever. Such a minimal change that could easily help the user keep
track of the currently focused scope/context and the actions they are
can be perform at this level.

\end{document}
